% =====================================================================
% Kapitel 3: Infrastruktur
% =====================================================================

\chapter{Infrastruktur}

\label{ch:infrastruktur}

Dieses Kapitel zeigt, wie alle Komponenten -- Datenbank, Python,
TypeScript, Tools -- auf einer Maschine zusammenstarten und
miteinander kommunizieren. Es ist die Grundlage, auf der alle
folgenden Kapitel aufbauen.

\section{Überblick: Was läuft wo?}

Abbildung~\ref{fig:infra-overview} zeigt, welche Prozesse auf der
Entwicklungsmaschine laufen und wie sie miteinander sprechen.

\begin{figure}[ht]
\centering
\begin{tikzpicture}[
  node distance=1cm and 1.5cm,
  box/.style={rectangle, draw, thick, rounded corners, minimum width=3cm, minimum height=1cm, align=center, font=\small},
  container/.style={rectangle, draw, very thick, dashed, rounded corners, inner sep=0.4cm},
  arrow/.style={-Latex, thick},
  dasharrow/.style={-Latex, thick, dashed},
]
  \node[container, label={\textbf{deephold\_db -- Entwicklungsmaschine}}] (host) {
    \begin{tikzpicture}[node distance=0.6cm]
      \node[box, fill=blue!20] (py) {Python venv\\\scriptsize deephold\_db package};
      \node[box, fill=green!20, right=of py] (pg) {PostgreSQL 16\\\scriptsize in Docker};
      \node[box, fill=yellow!20, right=of pg] (bun) {Bun\\\scriptsize deephold-db-tui};
      \node[box, fill=cyan!20, below=of pg] (admin) {Adminer\\\scriptsize Port 8080};
      \node[box, fill=orange!20, right=of admin] (prefect) {Prefect Server\\\scriptsize Port 4200};
    \end{tikzpicture}
  };

  \draw[arrow] (py) -- node[above, font=\scriptsize] {psycopg / SQLAlchemy} (pg);
  \draw[arrow] (bun) -- node[above, font=\scriptsize] {node-postgres} (pg);
  \draw[arrow] (pg) -- (admin);
  \draw[arrow] (pg) -- (prefect);
\end{tikzpicture}
\caption{Lokale Infrastruktur: Vier Prozesse, eine Datenbank. Alles
  läuft auf einer Maschine. PostgreSQL ist der zentrale Hub; Python
  und Bun kommunizieren nicht direkt, sondern nur über die DB.}
\label{fig:infra-overview}
\end{figure}

\begin{definition}
\textbf{localhost / 127.0.0.1}
\emph{Localhost} ist der Rechnername, der immer auf den
\emph{aktuellen Rechner} selbst verweist. Die IP-Adresse~127.0.0.1
ist die \emph{Loopback}-Adresse. Wenn ein Prozess auf Localhost hört,
ist er nur für Prozesse auf derselben Maschine erreichbar -- nicht
für andere Computer im Netz. Dies ist die sicherste
\emph{Standardtopologie} für Entwicklung.
\end{definition}

\section{Docker und Docker Compose}

\begin{definition}
\textbf{docker-compose.yml}
Die Datei \code{docker-compose.yml} ist die
\emph{Single Source of Truth} für alle containerisierten Services.
Sie listet jeden Service mit Image, Umgebungsvariablen, Ports,
Volumes und Healthcheck auf.
\end{definition}

\begin{lstlisting}[language=yaml, caption={Vollständige \code{docker-compose.yml}},label=lst:dockercompose]
name: deephold-db

services:
  postgres:
    image: postgres:16-alpine
    container_name: deephold_pg
    restart: unless-stopped
    environment:
      POSTGRES_USER: ${POSTGRES_USER:-deephold}
      POSTGRES_PASSWORD: ${POSTGRES_PASSWORD:-deephold}
      POSTGRES_DB: ${POSTGRES_DB:-deephold}
    ports:
      - "${POSTGRES_PORT:-5432}:5432"
    volumes:
      - postgres_data:/var/lib/postgresql/data
    healthcheck:
      test: ["CMD-SHELL", "pg_isready -U ${POSTGRES_USER:-deephold} -d ${POSTGRES_DB:-deephold}"]
      interval: 5s
      timeout: 5s
      retries: 10

  adminer:
    image: adminer:4
    container_name: deephold_adminer
    restart: unless-stopped
    ports:
      - "${ADMINER_PORT:-8080}:8080"
    environment:
      ADMINER_DEFAULT_SERVER: postgres
    depends_on:
      postgres:
        condition: service_healthy

  prefect-server:
    image: prefecthq/prefect:2.16-python3.11
    container_name: deephold_prefect
    restart: unless-stopped
    command: prefect server start --host 0.0.0.0
    ports:
      - "${PREFECT_PORT:-4200}:4200"
    environment:
      PREFECT_HOME: /opt/prefect
      PREFECT_SERVER_API_URL: http://localhost:4200/api
    volumes:
      - prefect_data:/opt/prefect

volumes:
  postgres_data:
  prefect_data:
\end{lstlisting}

Listing~\ref{lst:dockercompose} zeigt die vollständige
\code{docker-compose.yml}. Bemerkenswert:

\begin{itemize}
  \item \textbf{Drei Services}: \code{postgres} (Datenbank),
        \code{adminer} (Web-UI für SQL), \code{prefect-server}
        (Workflow-Orchestrierung).
  \item \textbf{Zwei Volumes}: \code{postgres\_data} und
        \code{prefect\_data} für persistente Daten.
  \item \textbf{Healthcheck} auf Postgres -- der \code{adminer}-Service
        startet erst, wenn Postgres bereit ist (\code{condition:
        service\_healthy}).
  \item \textbf{Umgebungsvariablen} mit Defaults:
        \code{\$\{POSTGRES\_USER:-deephold\}} -- falls die Variable
        nicht gesetzt ist, wird der Default \code{deephold} verwendet.
  \item \textbf{Image-Tags}: \code{postgres:16-alpine} -- die
        \code{-alpine}-Variante nutzt eine minimale Linux-Distribution
        (Alpine) und ist entsprechend klein.
\end{itemize}

\section{Das Makefile}

Das Makefile ist die \emph{einzige CLI}, die man kennen muss. Es
abstrahiert über alle Details -- Docker, Python, Bun, Alembic,
Tests.

\begin{lstlisting}[language=make, caption={Vollständiges Makefile},label=lst:makefile]
.PHONY: help
help:  ## Zeige alle Targets
	@grep -E '^[a-zA-Z_-]+:.*?## .*$$' $(MAKEFILE_LIST) | \
	  awk 'BEGIN {FS = ":.*?## "}; {printf "\033[36m%-20s\033[0m %s\n", $$1, $$2}'

.PHONY: init
init:  ## Komplett-Setup: deps + Stack + Migrations
	poetry install
	$(COMPOSE) up -d
	@echo "Warte auf Postgres ..."
	@for i in $$(seq 1 30); do \
	  $(COMPOSE) exec -T postgres pg_isready -U $${POSTGRES_USER:-deephold} >/dev/null 2>&1 && break; \
	  sleep 1; \
	done
	$(MAKE) migrate
	@echo "Init fertig. UI: Adminer :8080, Prefect :4200"

.PHONY: up
up:  ## docker-compose Stack starten
	$(COMPOSE) up -d

.PHONY: down
down:  ## docker-compose Stack stoppen
	$(COMPOSE) down

.PHONY: migrate
migrate:  ## Alembic-Migrationen anwenden
	DATABASE_URL=$(DB_DSN) $(PYTHON) -m alembic upgrade head

.PHONY: revision
revision:  ## Neue Alembic-Revision (Usage: make revision m="...")
	DATABASE_URL=$(DB_DSN) $(PYTHON) -m alembic revision --autogenerate -m "$(m)"

.PHONY: test
test:  ## pytest + pandera schemas
	DATABASE_URL=$(DB_DSN) $(PYTHON) -m pytest

.PHONY: tui
tui:  ## OpenTUI-Explorer starten (Bun erforderlich)
	@if ! command -v bun >/dev/null 2>&1; then \
		echo "Bun nicht installiert. Install: curl -fsSL https://bun.sh/install | bash"; \
		exit 1; \
	fi
	cd tui && bun install --silent && bun run start

.PHONY: tui-test
tui-test:  ## TUI-Tests (Unit + Integration)
	@if ! command -v bun >/dev/null 2>&1; then \
		echo "Bun nicht installiert. Install: curl -fsSL https://bun.sh/install | bash"; \
		exit 1; \
	fi
	cd tui && bun test

.PHONY: format
format:  ## ruff format
	$(PYTHON) -m ruff format .

.PHONY: lint
lint:  ## ruff check
	$(PYTHON) -m ruff check .

.PHONY: typecheck
typecheck:  ## mypy
	$(PYTHON) -m mypy src/

.PHONY: check
check:  ## Alles: lint + typecheck + test
	$(PYTHON) -m ruff check .
	$(PYTHON) -m ruff format --check .
	$(PYTHON) -m mypy src/
	$(PYTHON) -m pytest
\end{lstlisting}

Listing~\ref{lst:makefile} zeigt die wichtigsten Targets. Die
Konvention ist: \code{make <target>} -- und das Makefile sucht
sich die richtigen Befehle zusammen.

\section{Python-Environment}

\begin{definition}
\textbf{venv (Virtual Environment)}
Ein \emph{virtuelles Python-Environment} ist ein isoliertes
Verzeichnis, in dem Python-Pakete installiert werden, ohne das
System-Python zu beeinflussen. In diesem Projekt wird statt
Poetry (das viele zusätzliche Komplexität mitbringt) direkt
\code{python -m venv .venv} und \code{pip install} verwendet -- so
funktioniert das Setup auf jeder Maschine ohne spezielle Tools.
\end{definition}

\begin{beispiel}
\textbf{Setup-Sequenz}
\begin{lstlisting}[language=shell]
# 1. Virtuelles Environment anlegen
python3 -m venv .venv

# 2. pip aktualisieren
.venv/bin/pip install --quiet --upgrade pip

# 3. Projekt installieren (inkl. dev-deps)
.venv/bin/pip install --quiet -e .

# 4. Tests laufen lassen
.venv/bin/pytest
\end{lstlisting}
\end{beispiel}

Die wichtigsten Python-Pakete sind in
Tabelle~\ref{tab:py-deps} zusammengefasst.

\begin{table}[ht]
\centering
\small
\begin{tabular}{lll}
  \toprule
  \textbf{Paket} & \textbf{Version} & \textbf{Zweck} \\
  \midrule
  \code{polars}    & $\geq$0.20 & Datenverarbeitung (schneller als Pandas) \\
  \code{pandas}    & $\geq$2.2  & Datenverarbeitung (Kompatibilität) \\
  \code{numpy}     & $\geq$1.26 & Numerik \\
  \code{sqlalchemy} & $\geq$2.0  & ORM \\
  \code{alembic}   & $\geq$1.13 & Schema-Migrationen \\
  \code{psycopg[binary]} & $\geq$3.1 & PostgreSQL-Treiber \\
  \code{httpx}      & $\geq$0.27 & HTTP-Client (Vendor-Adapter) \\
  \code{tenacity}   & $\geq$8.2  & Retries \\
  \code{structlog}  & $\geq$24.1 & Logging \\
  \code{pydantic}   & $\geq$2.6  & Datenmodelle \\
  \code{pydantic-settings} & $\geq$2.2 & Config aus .env \\
  \code{prefect}    & $\geq$2.16 & Workflow-Orchestrierung \\
  \code{tabulate}   & $\geq$0.9  & Formatierte Tabellen im Terminal \\
  \midrule
  \code{pytest}     & $\geq$8.0  & Test-Runner \\
  \code{ruff}       & $\geq$0.3  & Linter + Formatter \\
  \code{mypy}       & $\geq$1.8  & Type-Checker \\
  \bottomrule
\end{tabular}
\caption{Wichtigste Python-Abhängigkeiten}
\label{tab:py-deps}
\end{table}

\section{Bun und TypeScript}

Bun ist die bevorzugte Runtime für das TUI (Kapitel~\ref{ch:tui}).
Es vereint Paketmanager, Test-Runner und Bundler in einem einzigen,
sehr schnellen Binary.

\begin{lstlisting}[language=shell, caption={Bun-Setup}]
# Installation (einmalig pro Maschine)
curl -fsSL https://bun.sh/install | bash
source ~/.bashrc

# Im tui/-Verzeichnis
cd tui
bun install                  # Installiert Dependencies
bun run start                # Startet TUI
bun test                     # Lässt alle Tests laufen
bun run typecheck            # tsc --noEmit
\end{lstlisting}

\section{Pfade, Ports, Datenverzeichnisse}

\begin{table}[ht]
\centering
\small
\begin{tabular}{lll}
  \toprule
  \textbf{Was} & \textbf{Pfad / Port} & \textbf{Wer greift zu} \\
  \midrule
  Postgres-Server         & \code{localhost:5432}    & Python, Bun, Adminer, psql \\
  Adminer-Web-UI          & \code{localhost:8080}    & Browser \\
  Prefect-Server-UI       & \code{localhost:4200}    & Browser \\
  PostgreSQL-Daten-Volume & \code{postgres\_data}    & Docker (persistent) \\
  Prefect-Daten-Volume    & \code{prefect\_data}     & Docker (persistent) \\
  Python-Pakete           & \code{./.venv/}          & \code{./.venv/bin/python} \\
  Bun-Cache               & \code{$\sim$/.bun/}      & Bun \\
  TUI-Node-Modules        & \code{./tui/node\_modules/} & Bun \\
  Jupyter-Notebook-Output & \code{./notebooks/}      & Browser (HTML) \\
  \bottomrule
\end{tabular}
\caption{Pfade und Ports aller Komponenten}
\label{tab:paths-ports}
\end{table}

\section{Was beim ersten Start passiert}

\begin{example}[Boot-Sequenz]
\begin{enumerate}
  \item \code{make up} startet die Docker-Container.
  \item Postgres initialisiert sich (5--10 Sekunden) und ist bereit
        für Verbindungen (\code{pg\_isready} gibt Exit-Code 0).
  \item Adminer wartet auf \code{postgres: service\_healthy},
        startet dann.
  \item Prefect-Server startet (langsamer, 10--20 Sekunden).
  \item \code{make migrate} führt alle Alembic-Migrationen aus und
        erzeugt die 14~Tabellen.
  \item \code{make ingest-<asset\_class>} oder \code{pytest} oder
        \code{make tui} können jetzt starten.
\end{enumerate}
\end{example}

\section{Was als Nächstes?}

Mit der Infrastruktur im Griff wenden wir uns dem nächsten
zentralen Thema zu: dem \emph{Datenmodell}. Kapitel~\ref{ch:datenmodell}
erklärt die 14~Tabellen, ihre Beziehungen und ihre Indizes.
